\documentclass[10pt,a4paper]{amsart}
\usepackage[margin=19mm]{geometry}
\usepackage{amsmath,amssymb,mathtools}
\usepackage[scr=rsfs,cal=euler]{mathalfa}
\usepackage{xfrac}
\usepackage[unicode,hidelinks,bookmarks=false]{hyperref}
\newcommand{\ie}{i.e.\ }
\newcommand{\cf}{cf.\ }
\newcommand{\resp}{respectively\ }
\newcommand{\Subsection}{\subsection}
\providecommand{\nobreakdash}{\nobreak\hbox{-}}
\setlength{\emergencystretch}{2em}
\multlinegap0pt
\usepackage{amssymb}
\usepackage{csquotes}
\usepackage{dsfont}
\usepackage{bm}
\usepackage{enumerate}
\usepackage{mathtools}
\newtheorem{theorem}{Theorem}
\newtheorem{lemma}{Lemma}
\title{A symmetric multivariate Elekes--R\'{o}nyai theorem}
\author{Yewen Sun}
\date{Source: arXiv:2504.02075v2 (CC BY 4.0)}
\begin{document}
\maketitle

\noindent\emph{Source and licence.} A symmetric multivariate Elekes--R\'{o}nyai theorem. Version arXiv:2504.02075v2 (CC BY 4.0), distributed by arXiv under the Creative Commons Attribution 4.0 International licence, http://creativecommons.org/licenses/by/4.0/, as recorded in the arXiv licence metadata for this exact version and shown on the abstract page https://arxiv.org/abs/2504.02075v2. Full licence notice: \texttt{licenses/arxiv-2504.02075-CC-BY-4.0.txt}.\par\smallskip

\noindent\textit{Editorial scope.} Counted unit: the article's theorem thm:main1 (source0.tex lines 121--139). The article's complete original proof of it is delivered with it (source0.tex lines 1330--1495, one \texttt{proof} environment of 166 lines, which the article places away from the statement). No mathematics is written, repaired or re-derived: the statement and the proof are the article's own lines, in the article's own notation, and no retained line is substituted or re-flowed.  Retained uncounted as the source-local context the proof needs: lines 171--190; lines 1165--1175; lines 1249--1261.  Nothing was omitted from any retained span, so the package carries no omission record.  No retained line contains a citation, so the wrapper carries no bibliography and no bibliography entry.  No retained line contains a figure environment, an image-inclusion command or drawing code, so no image is delivered and the package declares no asset dependency.  Editorial formatting only, in addition: an AMS \texttt{amsart} preamble that restates the article's own declarations and macros used by the retained lines, namely the environments \texttt{theorem}, \texttt{lemma} on separate counters, together with the standard packages those lines need.  The retained environments print in this document as Theorem 1, Lemma 1 in order of appearance, whereas the article prints the same items with the numbering of its own full document.  The complete native source of the article as distributed in the arXiv e-print for this exact version is bundled unchanged as \texttt{sources/175866/source0.tex}.
\begin{theorem}\label{thm:main2} Let $P(x_{1},\ldots,x_{d}),Q(x_{1},\ldots,x_{d})$ be two polynomials in $ \mathbb{R}[x_{1},\ldots,x_{d}]$  for some $d\geq 2.$ Assume $P,Q$ have degree at most $\delta$ and $P,Q$ depend  non-trivially on each of $x_1,...,x_d$. Let $1\leq t\leq d-1$ be a positive integer. Then for all $n \in \mathbb{N}$ and subsets $A_{1},\ldots, A_{d}$ of $\mathbb{R}$ satisfying \[|A_{i}|=n \text{~for each~} i\in [d],\] we have
$$
\max \{|P(A_{1},\ldots, A_{d})|,|Q(A_{1},\ldots, A_{d})|\} \gg_{\delta} n^{\frac{3}{2}-\frac{1}{2^{t+1}}},
$$
unless one of the following holds:

(i) $P$ and $Q$ form a $t$-additive pair, i.e., \begin{align*}&P(x_{1},\ldots,x_{d})=f( u_{1}(x_{1})+\cdots+ u_{d}(x_{d}))\quad \text{and}\\&Q(x_{1},\ldots,x_{d})=g( v_{1}(x_{1})+\cdots+v_{d}(x_{d})),\end{align*} where $f(x), g(x)$,$u_{1}(x),\ldots,u_{d}(x)$,$v_{1}(x),\ldots,v_{d}(x)$ are nonconstant univariate polynomials over $\mathbb{R}$. Moreover, \(\bigl|\{\,i\in[d]:\ u_i \not\equiv_a v_i\,\}\bigr|<t.\)






(ii) $P$ and $Q$ form a $t$-multiplicative pair, i.e.,  \begin{align*}&P(x_{1},\ldots,x_{d})=f( u_{1}(x_{1})\cdot\ldots\cdot u_{d}(x_{d}))\quad \text{and}\\&Q(x_{1},\ldots,x_{d})=g( v_{1}(x_{1}) \cdot\ldots\cdot v_{d}(x_{d})),\end{align*}
where $f(x), g(x)$,$u_{1}(x),\ldots,u_{d}(x)$,$v_{1}(x),\ldots,v_{d}(x)$  are nonconstant univariate polynomials over $\mathbb{R}$.  Moreover, \(\bigl|\{\,i\in[d]:\ u_i \not\equiv_m v_i\,\}\bigr|<t.\)




\end{theorem}

\begin{lemma}\label{lem:poly}
Let $d\geq 2$ and $f,g,u_{1},\ldots,u_{d},v_{1},\ldots,v_{d}\in \mathbb{R}[x]$ be nonconstant polynomials.

(i)  Assume $u_{i}$ and $v_{i}$ have no constants for all $i\in[d]$. If
$$f(u_{1}(x_{1})+\cdots+u_{d}(x_{d}))=g(v_{1}(x_{1})+\cdots+v_{d}(x_{d}))$$
for all $x_{1},\ldots,x_{d}\in\mathbb{R}$, then $u_{i}\equiv_{a} v_{i}$ for all $i\in[d]$.

(ii) Assume $u_{i}$ and $v_{i}$ are monic. If
$$f(u_{1}(x_{1})\cdot\ldots\cdot u_{d}(x_{d}))=g(v_{1}(x_{1}) \cdot\ldots\cdot v_{d}(x_{d}))$$
for all $x_{1},\ldots,x_{d}\in\mathbb{R}$, then $u_{i}\equiv_{m} v_{i}$ for all $i\in[d].$
\end{lemma}

\begin{lemma}\label{lem:counting}
Let $d\geq 1$ be a positive integer and let $\equiv$ be an equivalence relation on the set $[d] := \{1,2,\dots,d\}$.
Denote by $\mathfrak{S}_d$ the group of permutations on $[d]$.
For an integer $t\in [d]$, assume that for every permutation $\sigma \in \mathfrak{S}_d$, there exists a subset
$I_{\sigma}\subseteq [d]$ with $|I_{\sigma}|=t$ satisfying
\[
i\equiv \sigma(i) \quad \text{for all } i\in I_{\sigma}.
\]
Then there exists at least one equivalence class whose size is at least
\[
\left\lceil \frac{d+t}{2} \right\rceil.
\]
\end{lemma}

\begin{theorem}\label{thm:main1} 

Let $P(x_{1},\ldots,x_{d})$ be a polynomial in $ \mathbb{R}[x_{1},\ldots,x_{d}]$  for some $d\geq 2.$ Assume $P$ has degree $ \delta $ and $P$ depends non-trivially on each of $x_1,...,x_d$.  For any integer $t$ with $1\leq t\leq d-1$, $n \in \mathbb{N}$ and finite $A\subset \mathbb{R}$ with $|A|=n$, we have
$$
| P(A,\ldots, A)|\gg_{\delta} n^{\frac{3}{2}-\frac{1}{2^{t+1}}},
$$
unless one of the following holds:

(i) $P(x_{1},\ldots,x_{d})=f( u_{1}(x_{1})+\cdots +u_{d}(x_{d}))$ for some nonconstant univariate real polynomials  $f(x)$,$u_1(x),\ldots,u_{d}(x)$. 
Moreover, there exists an index subset $I\subseteq [d]$ with $|I|=d-\lfloor \frac{t-1}{2} \rfloor$ such that for any $i,j\in I$,  $u_{i}(x)\equiv_{a} u_{j}(x)$.



(ii) $P(x_{1},\ldots,x_{d})=f( u_{1}(x_{1})\cdot\ldots\cdot u_{d}(x_{d}))$ for some nonconstant univariate real polynomials  $f(x)$,$u_1(x),\ldots,u_{d}(x)$. Moreover, there exists an index subset $I\subseteq [d]$ with $|I|=d-\lfloor \frac{t-1}{2} \rfloor$ such that for any $i,j\in I$, $u_{i}(x)\equiv_{m} u_{j}(x)$.




\end{theorem}

\begin{proof}
Assume for contradiction that there exist $n\in\mathbb N$ and $A\subset\mathbb R$ with $|A|=n$ such that
$$|P(A,\ldots,A)|\ll_{\delta} n^{\frac32-\frac1{2^{t+1}}},$$
and that neither alternative (i) nor (ii) of Theorem~\ref{thm:main1} holds.

For $\sigma\in\mathfrak S_d$, write
$P^{\sigma}(x_1,\ldots,x_d):=P(x_{\sigma(1)},\ldots,x_{\sigma(d)})$.
Then $P^{\sigma}(A,\ldots,A)=P(A,\ldots,A)$ for every $\sigma$.
Applying Theorem~\ref{thm:main2} to the pair $(P,P^{\sigma})$ with $A_1=\cdots=A_d=A$, we have that for every $\sigma$
the pair $(P,P^{\sigma})$ forms either an additive pair or a multiplicative pair (since the expansion conclusion fails by assumption).

We first note that $P$ cannot be simultaneously of additive and multiplicative type.
Indeed, suppose for contradiction that
\[
P=f\big(u_1(x_1)+\cdots+u_d(x_d)\big)=g\big(w_1(x_1)\cdots w_d(x_d)\big)
\]
with all $u_i,w_i$ nonconstant.
Then in the fraction field of $\mathbb R[x_1,\ldots,x_d]$, for $i\neq j$ we have
\[
\frac{\partial_{x_i}P}{\partial_{x_j}P}
=\frac{u_i'(x_i)}{u_j'(x_j)}
=\frac{w_i'(x_i)/w_i(x_i)}{w_j'(x_j)/w_j(x_j)}.
\]
Fix $j$ and vary $x_i,x_j$; separation of variables implies there exists $c_i\in\mathbb R^\times$ such that
\[
u_i'(x)=c_i\,\frac{w_i'(x)}{w_i(x)}\qquad\text{in }\mathbb R(x).
\]
Clearing denominators gives $w_i(x)\,u_i'(x)=c_i\,w_i'(x)$ in $\mathbb R[x]$, hence $w_i\mid w_i'$.
This is impossible for a nonconstant polynomial (since $\deg w_i'<\deg w_i$), a contradiction.
Therefore, exactly one of the two types works for $P$.

\smallskip
\noindent\emph{Additive case.}
Assume $P$ is of additive type. Then (taking $\sigma=\mathrm{id}$ above) we may write
$P(x_1,\ldots,x_d)=f(u_1(x_1)+\cdots+u_d(x_d))$
with $f,u_1,\ldots,u_d$ nonconstant.
By subtracting constants from each $u_i$ and absorbing the total constant into $f$,
we may assume $u_i(0)=0$ for all $i$.


When we apply Theorem~\ref{thm:main2} to $(P,P^\sigma)$, the $t$-additive pair conclusion a priori provides
some representation $P=\widetilde f(\sum_{i=1}^d \widetilde u_i(x_i))$.
After shifting constants so that $\widetilde u_i(0)=0$ for all $i$ (absorbing the total constant into $\widetilde f$),
Lemma~\ref{lem:poly}(i) applied to
\[
f\Big(\sum_{i=1}^d u_i(x_i)\Big)=\widetilde f\Big(\sum_{i=1}^d \widetilde u_i(x_i)\Big)
\]
implies $\widetilde u_i\equiv_a u_i$ for every $i$.
Hence, up to $\equiv_a$, we may assume the $t$-additive pair decomposition of $P$ uses our fixed $u_i$.




Define an equivalence relation on $[d]$ by $i\sim j$ iff $u_i\equiv_a u_j$.

Fix an arbitrary $\sigma\in\mathfrak S_d$.
Since $P$ is not of multiplicative type, the pair $(P,P^{\sigma})$ must form a $t$-additive pair.
Thus there exist nonconstant univariate polynomials $g,v_1,\ldots,v_d$ such that
$$P^{\sigma}(x_1,\ldots,x_d)=g(v_1(x_1)+\cdots+v_d(x_d)),$$
and the mismatch set
$M_{\sigma}:=\{i\in[d]: u_i\not\equiv_a v_i\}$
satisfies $|M_{\sigma}|<t$.
On the other hand,
\begin{align*}
P^{\sigma}(x_1,\ldots,x_d)
&:=P(x_{\sigma(1)},\ldots,x_{\sigma(d)})\\
&=f\big(u_1(x_{\sigma(1)})+\cdots+u_d(x_{\sigma(d)})\big)\\
&=f\big(u_{\sigma^{-1}(1)}(x_1)+\cdots+u_{\sigma^{-1}(d)}(x_d)\big).
\end{align*}
After replacing each $v_i(x)$ by $v_i(x)-v_i(0)$ and composing $g$ with a translation (so that the equality defining $P^\sigma$ is unchanged),
we may assume $v_i(0)=0$ for all $i$; this normalization does not increase the mismatch set because $u_i(0)=0$ for all $i$, and if $u_i\equiv_a v_i$ then necessarily $v_i(0)=0$ as well.
Applying Lemma~\ref{lem:poly}(i) to
$$f(u_{\sigma^{-1}(1)}(x_1)+\cdots+u_{\sigma^{-1}(d)}(x_d))
=g(v_1(x_1)+\cdots+v_d(x_d))$$
gives $u_{\sigma^{-1}(i)}\equiv_a v_i$ for all $i\in[d]$.

Therefore, for every $\sigma$ there exists a set $I_{\sigma}\subseteq[d]$ with
$|I_{\sigma}|\ge d-(t-1)=d-t+1$
such that $i\sim \sigma^{-1}(i)$ for all $i\in I_{\sigma}$.
Since $\sigma$ ranges over all permutations, so does $\sigma^{-1}$, hence we may reindex and conclude:
for every $\tau\in\mathfrak S_d$ there exists $J_{\tau}\subseteq[d]$ with $|J_{\tau}|=d-t+1$ such that
$i\sim \tau(i)$ for all $i\in J_{\tau}$.

Now apply Lemma~\ref{lem:counting} with this equivalence relation and with the parameter
$t':=d-t+1$.
We obtain an equivalence class $S\subseteq[d]$ with
$|S|\ge \left\lceil \frac{d+t'}2\right\rceil
=\left\lceil \frac{2d-t+1}{2}\right\rceil
=d-\left\lfloor \frac{t-1}{2}\right\rfloor$,
and $u_i\equiv_a u_j$ for all $i,j\in S$.
This is exactly alternative (i) in Theorem~\ref{thm:main1} (taking $I=S$).

\smallskip
\noindent\emph{Multiplicative case.}
Assume $P$ is of multiplicative type. Then (taking $\sigma=\mathrm{id}$ above) we may write
\[
P(x_1,\ldots,x_d)=f\big(u_1(x_1)\cdots u_d(x_d)\big),
\]
with $f,u_1,\ldots,u_d$ nonconstant.
By dividing each $u_i$ by its leading coefficient and absorbing the resulting nonzero constant into $f$,
we may assume that every $u_i$ is monic.

Similarly, in the multiplicative setting, any representation $P=\widetilde f(\prod_{i=1}^d \widetilde u_i(x_i))$
can be normalized so that each $\widetilde u_i$ is monic (absorbing constants into $\widetilde f$).
Then Lemma~\ref{lem:poly}(ii) implies $\widetilde u_i\equiv_m u_i$ for all $i$.
Thus we may assume the $t$-multiplicative pair decomposition of $P$ uses our fixed monic $u_i$, up to $\equiv_m$.



Define an equivalence relation on $[d]$ by $i\sim j$ iff $u_i\equiv_m u_j$.

Fix an arbitrary $\sigma\in\mathfrak S_d$.
Since $P$ is not of additive type, the pair $(P,P^\sigma)$ must form a $t$-multiplicative pair.
Thus there exist nonconstant univariate polynomials $g,v_1,\ldots,v_d$ such that
\[
P^{\sigma}(x_1,\ldots,x_d)=g\big(v_1(x_1)\cdots v_d(x_d)\big),
\]
and the mismatch set
\[
M_{\sigma}:=\{i\in[d]: u_i\not\equiv_m v_i\}
\]
satisfies $|M_\sigma|<t$.
Normalizing as above, we may assume each $v_i$ is monic: write $v_i(x)=c_i\,\widetilde v_i(x)$ with $\widetilde v_i$ monic and $c_i\in\mathbb R^\times$, and absorb $\prod_i c_i$ into the outer polynomial $g$.
This does not increase the mismatch set, because if $u_i$ is monic and $u_i\equiv_m v_i$ then the leading coefficient of $v_i$ has $|c_i|=1$, and dividing by such a constant does not affect $\equiv_m$.


On the other hand,
\begin{align*}
P^{\sigma}(x_1,\ldots,x_d)
&=P(x_{\sigma(1)},\ldots,x_{\sigma(d)})\\
&=f\big(u_1(x_{\sigma(1)})\cdots u_d(x_{\sigma(d)})\big)\\
&=f\big(u_{\sigma^{-1}(1)}(x_1)\cdots u_{\sigma^{-1}(d)}(x_d)\big).
\end{align*}
Applying Lemma~\ref{lem:poly}(ii) to the identity
\[
f\big(u_{\sigma^{-1}(1)}(x_1)\cdots u_{\sigma^{-1}(d)}(x_d)\big)
=
g\big(v_1(x_1)\cdots v_d(x_d)\big),
\]
we obtain
\[
u_{\sigma^{-1}(i)}\equiv_m v_i \qquad \text{for all } i\in[d].
\]
Hence for every $i\in[d]\setminus M_\sigma$ we have
\[
u_i\equiv_m v_i\equiv_m u_{\sigma^{-1}(i)},
\]
i.e.\ $i\sim \sigma^{-1}(i)$.

Therefore, for every $\sigma$ there exists a set $I_\sigma\subseteq[d]$ with
$|I_\sigma|\ge d-(t-1)=d-t+1$ such that $i\sim \sigma^{-1}(i)$ for all $i\in I_\sigma$.
Since $\sigma$ ranges over all permutations, so does $\sigma^{-1}$, hence:
for every $\tau\in\mathfrak S_d$ there exists $J_\tau\subseteq[d]$ with $|J_\tau|=d-t+1$ such that
$i\sim \tau(i)$ for all $i\in J_\tau$.

Now apply Lemma~\ref{lem:counting} with this equivalence relation and with parameter $t':=d-t+1$.
We obtain an equivalence class $S\subseteq[d]$ with
\[
|S|\ge \left\lceil \frac{d+t'}2\right\rceil
=\left\lceil \frac{2d-t+1}{2}\right\rceil
=d-\left\lfloor \frac{t-1}{2}\right\rfloor,
\]
and $u_i\equiv_m u_j$ for all $i,j\in S$.
This is exactly alternative (ii) in Theorem~\ref{thm:main1} (taking $I=S$).

\end{proof}

\end{document}
